\chapter{การประเมินผลและการปรับปรุงประสิทธิภาพของตัวแบบ}
\label{CH7_Evaluation}

หลังจากที่ได้ศึกษาวิธีการสร้างตัวแบบการเรียนรู้ของเครื่องหลายรูปแบบในบทก่อนหน้า ไม่ว่าจะเป็นการถดถอยเชิงเส้น การถดถอยโลจิสติค ต้นไม้ตัดสินใจ และการสุ่มป่าไม้ คำถามสำคัญที่ตามมาคือ ``ตัวแบบที่สร้างขึ้นมีประสิทธิภาพเพียงพอหรือไม่'' และ ``เราจะเปรียบเทียบตัวแบบหลายตัวแบบเพื่อเลือกตัวที่ดีที่สุดได้อย่างไร'' บทนี้จะอธิบายเครื่องมือและวิธีการทางสถิติที่จำเป็นสำหรับการประเมินประสิทธิภาพของตัวแบบอย่างเป็นระบบ โดยแบ่งเนื้อหาออกเป็น 5 ส่วนหลัก ได้แก่ การประเมินประสิทธิภาพสำหรับการพยากรณ์ (Regression) การประเมินประสิทธิภาพสำหรับการจำแนกประเภท (Classification) เส้นโค้ง ROC และค่า AUC การตรวจสอบความถูกต้องของตัวแบบด้วย Train-Test Split และ Cross-Validation และส่วนสุดท้ายครอบคลุมปัญหา Overfitting และ Underfitting รวมถึงเทคนิค Regularization

%##########################################
\section{การประเมินประสิทธิภาพสำหรับการพยากรณ์}
การประเมินประสิทธิภาพของตัวแบบการพยากรณ์ (Regression Models) เป็นขั้นตอนที่มีความสำคัญยิ่งในการวิเคราะห์ข้อมูล เนื่องจากตัวแบบการพยากรณ์ที่ดีต้องสามารถทำนายค่าที่ใกล้เคียงกับค่าจริงมากที่สุด ในทางคณิตศาสตร์ ความคลาดเคลื่อน (Error) ของการทำนายนิยามได้จากผลต่างระหว่างค่าจริง $y_i$ และค่าทำนาย $\hat{y}_i$ กล่าวคือ $e_i = y_i - \hat{y}_i$ และเป้าหมายของการเรียนรู้คือการหาตัวแบบที่ทำให้ความคลาดเคลื่อนเหล่านี้มีค่าน้อยที่สุดโดยรวม การวัดประสิทธิภาพจึงต้องอาศัยตัวชี้วัดที่สามารถรวมความคลาดเคลื่อนทั้งหมดในชุดข้อมูลให้เป็นค่าตัวเลขเดียวที่บ่งบอกคุณภาพของตัวแบบได้อย่างชัดเจน

\subsection{MSE, RMSE, MAE และ R-squared}
ในการวัดความคลาดเคลื่อนของตัวแบบการพยากรณ์ มีตัวชี้วัดมาตรฐานที่ใช้กันอย่างแพร่หลาย 4 ตัว ได้แก่ MSE (Mean Squared Error) RMSE (Root Mean Squared Error) MAE (Mean Absolute Error) และ R-squared ซึ่งแต่ละตัวมีความหมายและการตีความที่แตกต่างกัน

\begin{myBox}[]{สูตรคำนวณตัวชี้วัดประสิทธิภาพการพยากรณ์}{}
สมมติให้มีข้อมูล $n$ ตัวอย่าง โดยที่ $y_i$ คือค่าจริงและ $\hat{y}_i$ คือค่าทำนายของตัวแบบ ตัวชี้วัดทั้ง 4 คำนวณได้ดังนี้

\begin{enumerate}
\item \textbf{Mean Squared Error (MSE)}
$$MSE = \frac{1}{n}\sum_{i=1}^{n}(y_i - \hat{y}_i)^2 = \frac{1}{n}\sum_{i=1}^{n}e_i^2$$

\item \textbf{Root Mean Squared Error (RMSE)}
$$RMSE = \sqrt{MSE} = \sqrt{\frac{1}{n}\sum_{i=1}^{n}(y_i - \hat{y}_i)^2}$$

\item \textbf{Mean Absolute Error (MAE)}
$$MAE = \frac{1}{n}\sum_{i=1}^{n}|y_i - \hat{y}_i| = \frac{1}{n}\sum_{i=1}^{n}|e_i|$$

\item \textbf{R-squared ($R^2$)}
$$R^2 = 1 - \frac{\sum_{i=1}^{n}(y_i - \hat{y}_i)^2}{\sum_{i=1}^{n}(y_i - \bar{y})^2} = 1 - \frac{SS_{res}}{SS_{tot}}$$
โดยที่ $\bar{y} = \frac{1}{n}\sum_{i=1}^{n}y_i$ คือค่าเฉลี่ยของค่าจริง
\end{enumerate}
\end{myBox}

ความแตกต่างระหว่าง MSE และ MAE อยู่ที่การให้น้ำหนักกับค่าคลาดเคลื่อน โดย MSE ยกกำลังสองค่าความคลาดเคลื่อนทำให้ค่าที่คลาดเคลื่อนมากถูกลงโทษหนักกว่าค่าที่คลาดเคลื่อนน้อย ดังนั้น MSE จึงไวต่อค่าผิดปกติ (Outliers) มากกว่า MAE ในทางกลับกัน MAE ให้น้ำหนักเท่ากันกับทุกค่าคลาดเคลื่อน ทำให้เป็นตัวชี้วัดที่ทนทาน (Robust) ต่อค่าผิดปกติมากกว่า ส่วน RMSE เป็นรากที่สองของ MSE ซึ่งมีข้อดีคือมีหน่วยเดียวกับตัวแปรตาม ทำให้ตีความได้ง่ายกว่า MSE

สำหรับ R-squared เป็นตัวชี้วัดที่มีความพิเศษ ค่า $R^2$ บ่งบอกสัดส่วนของความแปรปรวน (Variance) ในตัวแปรตามที่ตัวแบบสามารถอธิบายได้ ค่า $R^2$ มีค่าอยู่ระหว่าง 0 ถึง 1 โดยค่าใกล้ 1 หมายถึงตัวแบบมีความสามารถในการอธิบายความแปรปรวนของข้อมูลสูง ส่วนค่าใกล้ 0 หมายถึงตัวแบบไม่สามารถอธิบายความแปรปรวนได้ดี อย่างไรก็ตาม ค่า $R^2$ ที่สูงมากผิดปกติอาจเป็นสัญญาณของ Overfitting ซึ่งจะกล่าวถึงในรายละเอียดในส่วนหลังของบทนี้

\begin{exmp}
พิจารณาข้อมูลการพยากรณ์ยอดขายรายวันของร้านค้าปลีกแห่งหนึ่ง สมมติว่ามีข้อมูล 5 วันดังแสดงในตาราง \ref{table:sales_error} โดย $y_i$ คือยอดขายจริง (พันบาท) และ $\hat{y}_i$ คือยอดขายที่ตัวแบบทำนาย

\begin{table}[h]
\centering
\caption{ข้อมูลยอดขายจริงเทียบกับค่าทำนายของตัวแบบ}
\label{table:sales_error}
\begin{tabular}{|c|c|c|c|}
\hline
\textbf{วัน} & \textbf{ยอดขายจริง $y_i$} & \textbf{ทำนาย $\hat{y}_i$} & \textbf{คลาดเคลื่อน $e_i$} \\ \hline
1 & 45 & 43 & 2 \\ \hline
2 & 52 & 50 & 2 \\ \hline
3 & 38 & 42 & -4 \\ \hline
4 & 61 & 58 & 3 \\ \hline
5 & 55 & 57 & -2 \\ \hline
\end{tabular}
\end{table}

จากข้อมูลในตาราง คำนวณหาค่าตัวชี้วัดต่าง ๆ ได้ดังนี้

การคำนวณค่าคลาดเคลื่อน $e_i = y_i - \hat{y}_i$ แสดงดังคอลัมน์สุดท้าย โดยผลรวมของค่าคลาดเคลื่อนเท่ากับ $2+2-4+3-2 = 1$ ซึ่งใกล้ศูนย์แสดงว่าตัวแบบไม่มีความลำเอียงเชิงระบบ

สำหรับ MSE
$$MSE = \frac{1}{5}(2^2 + 2^2 + (-4)^2 + 3^2 + (-2)^2) = \frac{1}{5}(4+4+16+9+4) = \frac{37}{5} = 7.4$$

สำหรับ RMSE
$$RMSE = \sqrt{7.4} \approx 2.72 \text{ (พันบาท)}$$

สำหรับ MAE
$$MAE = \frac{1}{5}(|2| + |2| + |-4| + |3| + |-2|) = \frac{1}{5}(2+2+4+3+2) = \frac{13}{5} = 2.6 \text{ (พันบาท)}$$

สมมติค่าเฉลี่ยยอดขายจริง $\bar{y} = \frac{45+52+38+61+55}{5} = \frac{251}{5} = 50.2$

$$R^2 = 1 - \frac{\sum(y_i - \hat{y}_i)^2}{\sum(y_i - \bar{y})^2} = 1 - \frac{37}{(45-50.2)^2+(52-50.2)^2+(38-50.2)^2+(61-50.2)^2+(55-50.2)^2}$$
$$= 1 - \frac{37}{27.04+3.24+148.84+116.64+23.04} = 1 - \frac{37}{318.8} \approx 0.884$$

ค่า $R^2 \approx 0.884$ หมายความว่าตัวแบบสามารถอธิบายความแปรปรวนของยอดขายได้ประมาณร้อยละ 88.4 ซึ่งถือว่ามีประสิทธิภาพดี

\begin{figure}[h]
\begin{pycode}{regression\_metrics.py}
import numpy as np
from sklearn.metrics import (
    mean_squared_error,
    mean_absolute_error,
    r2_score
)

y_true = np.array([45, 52, 38, 61, 55])
y_pred = np.array([43, 50, 42, 58, 57])

mse = mean_squared_error(y_true, y_pred)
rmse = np.sqrt(mse)
mae = mean_absolute_error(y_true, y_pred)
r2 = r2_score(y_true, y_pred)

print(f"MSE: {mse:.2f}")
print(f"RMSE: {rmse:.2f}")
print(f"MAE: {mae:.2f}")
print(f"R2: {r2:.4f}")
\tcblower
ผลลัพธ์ที่ได้:

MSE: 7.40
RMSE: 2.72
MAE: 2.60
R2: 0.8840
\end{pycode}
\caption{Python code และผลลัพธ์การคำนวณตัวชี้วัดประสิทธิภาพการพยากรณ์}
\label{fig:regression_metrics_code}
\end{figure}
\end{exmp}

%##########################################
\section{การประเมินประสิทธิภาพสำหรับการจำแนกประเภท}
การประเมินประสิทธิภาพของตัวแบบการจำแนกประเภท (Classification Models) มีความซับซ้อนและแตกต่างจากการพยากรณ์อย่างมีนัยสำคัญ เนื่องจากผลลัพธ์ของการจำแนกประเภทเป็น ``หมวดหมู่'' (Class) ไม่ใช่ค่าตัวเลขที่ต่อเนื่อง ดังนั้นการวัด ``ความคลาดเคลื่อน'' จึงต้องพิจารณาว่าตัวแบบจำแนกถูกต้องหรือผิดในแต่ละกรณี และประเภทของความผิดพลาดที่เกิดขึ้นมีผลกระทบแตกต่างกันในทางปฏิบัติ ตัวอย่างเช่น ในระบบวินิจฉัยโรค การทำนายผู้ป่วยว่า ``ไม่เป็นโรค'' เมื่อจริง ๆ แล้วเป็นโรค (False Negative) มีผลร้ายแรงกว่าการทำนายว่า ``เป็นโรค'' เมื่อจริง ๆ ไม่เป็น (False Positive) อย่างมาก ดังนั้นการเลือกตัวชี้วัดที่เหมาะสมจึงขึ้นอยู่กับบริบทการใช้งาน

\subsection{Confusion Matrix}
Confusion Matrix เป็นเครื่องมือพื้นฐานที่ใช้แสดงผลการทำนายของตัวแบบการจำแนกประเภทในรูปตาราง โดยเปรียบเทียบผลการทำนายกับค่าจริง สำหรับปัญหาการจำแนกแบบทวิภาค (Binary Classification) ที่มี 2 กลุ่ม คือ Positive (P) และ Negative (N) โครงสร้างของ Confusion Matrix แสดงดังตาราง \ref{table:confusion_matrix}

\begin{table}[h]
\centering
\caption{โครงสร้างของ Confusion Matrix สำหรับการจำแนกแบบทวิภาค}
\label{table:confusion_matrix}
\begin{tabular}{|c|c|c|c|}
\hline
 & \multicolumn{2}{c|}{\textbf{ทำนาย}} \\ \cline{2-3}
\textbf{ค่าจริง} & \textbf{Positive} & \textbf{Negative} \\ \hline
\textbf{Positive} & True Positive (TP) & False Negative (FN) \\ \hline
\textbf{Negative} & False Positive (FP) & True Negative (TN) \\ \hline
\end{tabular}
\end{table}

จากตารางข้างต้น คำอธิบายของแต่ละช่องมีดังนี้

\begin{itemize}
\item \textbf{True Positive (TP)} คือจำนวนกรณีที่ค่าจริงเป็น Positive และตัวแบบทำนายถูกต้องว่าเป็น Positive

\item \textbf{True Negative (TN)} คือจำนวนกรณีที่ค่าจริงเป็น Negative และตัวแบบทำนายถูกต้องว่าเป็น Negative

\item \textbf{False Positive (FP)} คือจำนวนกรณีที่ค่าจริงเป็น Negative แต่ตัวแบบทำนายผิดว่าเป็น Positive (เรียกว่า Type I Error)

\item \textbf{False Negative (FN)} คือจำนวนกรณีที่ค่าจริงเป็น Positive แต่ตัวแบบทำนายผิดว่าเป็น Negative (เรียกว่า Type II Error)
\end{itemize}

\begin{exmp}
พิจารณาระบบวินิจฉัยโรคเบาหวานโดยใช้ตัวแบบการเรียนรู้ของเครื่อง โดยผลการทดสอบกับผู้ป่วย 1000 คน แสดงผลดังนี้: มีผู้ป่วยจริง 200 คน (Positive) และไม่เป็นโรค 800 คน (Negative) เมื่อนำผลการวินิจฉัยจากตัวแบบมาเปรียบเทียบกับผลตรวจที่ถูกต้อง พบว่า TP = 160, FN = 40, FP = 60, TN = 740

จากข้อมูลนี้ สามารถคำนวณตัวชี้วัดต่าง ๆ ได้ดังนี้

Overall Accuracy = $\frac{TP + TN}{TP + TN + FP + FN} = \frac{160 + 740}{1000} = 0.90$ หรือร้อยละ 90

Sensitivity (Recall ของ positive class) = $\frac{TP}{TP + FN} = \frac{160}{160 + 40} = 0.80$ หมายความว่าตัวแบบสามารถตรวจพบผู้เป็นโรคได้ร้อยละ 80

Specificity = $\frac{TN}{TN + FP} = \frac{740}{740 + 60} = 0.925$ หมายความว่าตัวแบบสามารถระบุผู้ไม่เป็นโรคได้ถูกต้องร้อยละ 92.5
\end{exmp}

\begin{figure}[h]
\begin{pycode}{confusion\_matrix.py}
from sklearn.metrics import (
    confusion_matrix,
    accuracy_score,
    recall_score,
    precision_score
)
import numpy as np

y_true = np.array([1]*200 + [0]*800)
y_pred = np.array([1]*160 + [0]*40 +
                   [1]*60 + [0]*740)

cm = confusion_matrix(y_true, y_pred)
acc = accuracy_score(y_true, y_pred)
sens = recall_score(y_true, y_pred,
                    pos_label=1)
spec = recall_score(y_true, y_pred,
                    pos_label=0)

print(f"Confusion Matrix:\n{cm}")
print(f"Accuracy: {acc:.4f}")
print(f"Sensitivity: {sens:.4f}")
print(f"Specificity: {spec:.4f}")
\tcblower
ผลลัพธ์:

Confusion Matrix:
[[740  60]
 [ 40 160]]
Accuracy: 0.9000
Sensitivity: 0.8000
Specificity: 0.9250
\end{pycode}
\caption{Python code และผลลัพธ์การคำนวณ Confusion Matrix}
\label{fig:confusion_matrix_code}
\end{figure}

\subsection{ตัวชี้วัดประสิทธิภาพ: Accuracy, Precision, Recall และ F1-Score}
จากค่าใน Confusion Matrix สามารถคำนวณตัวชี้วัดประสิทธิภาพได้หลายตัว ซึ่งแต่ละตัวเน้นมุมมองที่แตกต่างกันในการประเมินคุณภาพของตัวแบบ

\begin{myBox}[]{ตัวชี้วัดประสิทธิภาพการจำแนกประเภท}{}
\begin{enumerate}
\item \textbf{Accuracy} หรือความแม่นยำรวม แสดงสัดส่วนของการทำนายที่ถูกต้องทั้งหมด
$$Accuracy = \frac{TP + TN}{TP + TN + FP + FN}$$

\item \textbf{Precision} หรือความแม่นยำเชิงทำนาย แสดงสัดส่วนของกรณีที่ทำนายว่าเป็น Positive แล้วถูกต้องจริง ๆ
$$Precision = \frac{TP}{TP + FP}$$

\item \textbf{Recall} หรือความไว (Sensitivity) แสดงสัดส่วนของกรณี Positive จริงที่ตัวแบบสามารถตรวจพบได้
$$Recall = \frac{TP}{TP + FN}$$

\item \textbf{F1-Score} คือค่าเฉลี่ยฮาร์มอนิก (Harmonic Mean) ของ Precision และ Recall ซึ่งเป็นตัวชี้วัดที่สมดุลระหว่างทั้งสองตัว
$$F1\text{-}Score = 2 \times \frac{Precision \times Recall}{Precision + Recall}$$
\end{enumerate}
\end{myBox}

ความสัมพันธ์ระหว่าง Precision และ Recall มีลักษณะเป็น Trade-off โดยทั่วไปเมื่อต้องการเพิ่ม Recall อาจต้องลด Threshold ของความน่าจะเป็นที่ใช้ตัดสินใจ ทำให้ตัวแบบทำนาย Positive มากขึ้นแต่อาจมี False Positive เพิ่มขึ้นด้วย ทำให้ Precision ลดลง F1-Score จึงเป็นตัวชี้วัดที่เหมาะสมเมื่อต้องการสมดุลระหว่างทั้งสองตัว

\begin{exmp}
ใช้ข้อมูลจากตัวอย่างระบบวินิจฉัยโรคเบาหวานข้างต้น โดย TP = 160, FN = 40, FP = 60, TN = 740

คำนวณ Precision = $\frac{160}{160 + 60} = \frac{160}{220} \approx 0.727$ 

คำนวณ Recall = $\frac{160}{160 + 40} = \frac{160}{200} = 0.80$ (ซึ่งเท่ากับ Sensitivity ในกรณีทวิภาค)

คำนวณ F1-Score = $2 \times \frac{0.727 \times 0.80}{0.727 + 0.80} = 2 \times \frac{0.5816}{1.527} \approx 0.762$

จากผลการวิเคราะห์ ค่า Accuracy เท่ากับ 0.90 แสดงว่าตัวแบบทำนายถูกต้องร้อยละ 90 แต่เมื่อพิจารณา Precision = 0.727 ซึ่งหมายความว่าเมื่อตัวแบบทำนายว่าผู้ป่วยเป็นโรคเบาหวาน จะถูกต้องประมาณร้อยละ 72.7 เท่านั้น ส่วน Recall = 0.80 แสดงว่าตัวแบบสามารถตรวจพบผู้เป็นโรคได้ร้อยละ 80 ของผู้ที่เป็นโรคจริง ๆ ค่า F1-Score = 0.762 เป็นค่าเฉลี่ยที่สมดุลระหว่าง Precision และ Recall
\end{exmp}

%##########################################
\section{ROC Curve และ AUC}
เส้นโค้ง ROC (Receiver Operating Characteristic Curve) และค่า AUC (Area Under the ROC Curve) เป็นเครื่องมือที่ใช้ประเมินประสิทธิภาพของตัวแบบการจำแนกประเภทในสถานการณ์ที่มี Trade-off ระหว่าง True Positive Rate (TPR) และ False Positive Rate (FPR) ที่แตกต่างกันตาม Threshold ของการตัดสินใจ

\begin{myBox}[]{องค์ประกอบของ ROC Curve}{}
สำหรับการจำแนกแบบทวิภาค องค์ประกอบหลักของ ROC Curve ประกอบด้วย

\begin{enumerate}
\item \textbf{True Positive Rate (TPR)} หรือ Sensitivity
$$TPR = \frac{TP}{TP + FN} = Recall$$

\item \textbf{False Positive Rate (FPR)} หรือ 1 - Specificity
$$FPR = \frac{FP}{FP + TN} = 1 - Specificity$$
\end{enumerate}

เมื่อเปลี่ยน Threshold ของการตัดสินใจจาก 0 ถึง 1 จะได้ค่า TPR และ FPR ที่แตกต่างกัน ซึ่งเมื่อนำมาลงจุดบนกราฟโดยมี FPR เป็นแกน X และ TPR เป็นแกน Y จะเกิดเส้นโค้ง ROC

\begin{itemize}
\item ถ้า Threshold = 0 จะทำนายทุกตัวอย่างว่าเป็น Positive ทำให้ TPR = 1 และ FPR = 1 (จุดบนขวาบน)
\item ถ้า Threshold = 1 จะทำนายทุกตัวอย่างว่าเป็น Negative ทำให้ TPR = 0 และ FPR = 0 (จุดบนซ้ายล่าง)
\item ตัวแบบที่ทำนายสุ่มจะมี ROC Curve เป็นเส้นทแยงมุม y = x
\end{itemize}
\end{myBox}

\begin{exmp}
พิจารณาตัวแบบการจำแนกประเภทผู้เป็นโรคหัวใจที่ให้ความน่าจะเป็นในการเป็นโรคสำหรับผู้ป่วย 8 คน ดังแสดงในตาราง \ref{table:roc_example}

\begin{table}[h]
\centering
\caption{ผลการทำนายความน่าจะเป็นของตัวแบบและค่าจริง}
\label{table:roc_example}
\begin{tabular}{|c|c|c|}
\hline
\textbf{ผู้ป่วย} & \textbf{ความน่าจะเป็น $\hat{p}$} & \textbf{ค่าจริง} \\ \hline
1 & 0.90 & Positive \\ \hline
2 & 0.75 & Positive \\ \hline
3 & 0.60 & Positive \\ \hline
4 & 0.55 & Negative \\ \hline
5 & 0.45 & Positive \\ \hline
6 & 0.35 & Negative \\ \hline
7 & 0.20 & Negative \\ \hline
8 & 0.10 & Negative \\ \hline
\end{tabular}
\end{table}

เมื่อใช้ Threshold = 0.5 เป็นเกณฑ์ตัดสิน ผู้ที่มี $\hat{p} \geq 0.5$ จะถูกทำนายว่า Positive ดังนั้นผู้ป่วย 1, 2, 3, 4, 5 จะถูกทำนายว่า Positive และผู้ป่วย 6, 7, 8 จะถูกทำนายว่า Negative

Confusion Matrix ที่ Threshold = 0.5: TP = 3 (ผู้ป่วย 1, 2, 3), FN = 1 (ผู้ป่วย 5), FP = 1 (ผู้ป่วย 4), TN = 3 (ผู้ป่วย 6, 7, 8)

ดังนั้น TPR = 3/(3+1) = 0.75 และ FPR = 1/(1+3) = 0.25

หากเปลี่ยน Threshold = 0.6 จะได้ TPR = 2/(2+2) = 0.50 และ FPR = 0/(0+4) = 0 ซึ่งให้จุดที่แตกต่างกันบน ROC Curve

เมื่อวาด ROC Curve จากทุก Threshold ที่เป็นไปได้ พื้นที่ใต้เส้นโค้ง (AUC) จะบอกถึงประสิทธิภาพโดยรวมของตัวแบบ โดย AUC = 1 หมายถึงตัวแบบที่สมบูรณ์แบบ AUC = 0.5 หมายถึงตัวแบบที่ทำนายสุ่ม และ AUC < 0.5 หมายถึงตัวแบบที่แย่กว่าการทำนายสุ่ม
\end{exmp}

\begin{figure}[h]
\begin{pycode}{roc\_curve.py}
import numpy as np
import matplotlib.pyplot as plt
from sklearn.metrics import (
    roc_curve, auc, roc_auc_score
)

y_true = np.array([1,1,1,0,1,0,0,0])
y_prob = np.array([0.90,0.75,0.60,0.55,
                   0.45,0.35,0.20,0.10])

fpr, tpr, thresholds = roc_curve(
    y_true, y_prob)
roc_auc = auc(fpr, tpr)

plt.figure(figsize=(8,6))
plt.plot(fpr, tpr,
         label=f'ROC curve (AUC = {roc_auc:.3f})')
plt.plot([0,1],[0,1],'k--',label='Random')
plt.xlabel('False Positive Rate')
plt.ylabel('True Positive Rate')
plt.title('ROC Curve')
plt.legend()
plt.grid()
plt.show()
print(f"AUC: {roc_auc:.4f}")
\tcblower
ผลลัพธ์:

AUC: 0.8750

ROC Curve แสดงว่าตัวแบบมีประสิทธิภาพดี
โดยมี AUC = 0.875 ซึ่งสูงกว่าเส้นทแยงมุม (0.5)
อย่างมีนัยสำคัญ
\end{pycode}
\caption{Python code และผลลัพธ์การวาด ROC Curve}
\label{fig:roc_code}
\end{figure}

%##########################################
\section{การตรวจสอบความถูกต้องของตัวแบบ}
การตรวจสอบความถูกต้อง (Validation) ของตัวแบบการเรียนรู้ของเครื่องเป็นกระบวนการที่ใช้ประเมินว่าตัวแบบที่สร้างขึ้นมีความสามารถในการทำนายข้อมูลที่ไม่เคยเห็นมาก่อนได้ดีเพียงใด ทั้งนี้เนื่องจากวัตถุประสงค์หลักของการเรียนรู้ของเครื่องคือการสร้างตัวแบบที่สามารถนำไปใช้งานจริงกับข้อมูลใหม่ที่ยังไม่เคยเห็น ไม่ใช่เพียงการจำข้อมูลฝึกสอน (Training Data) การตรวจสอบความถูกต้องจึงเป็นขั้นตอนที่ขาดไม่ได้ในกระบวนการพัฒนาตัวแบบ เพื่อให้มั่นใจว่าตัวแบบมีความสามารถในการสังเคราะห์ (Generalization) ไปยังข้อมูลใหม่ได้อย่างแม่นยำ

\subsection{Train-Test Split}
วิธีการพื้นฐานที่สุดในการตรวจสอบความถูกต้องของตัวแบบคือการแบ่งข้อมูลออกเป็นชุดฝึกสอน (Training Set) และชุดทดสอบ (Test Set) โดยใช้ข้อมูลชุดฝึกสอนในการเรียนรู้และปรับแต่งตัวแบบ จากนั้นจึงใช้ข้อมูลชุดทดสอบในการประเมินประสิทธิภาพที่แท้จริงของตัวแบบ วิธีนี้เรียกว่า Train-Test Split ซึ่งมีความเรียบง่ายและเข้าใจได้ง่าย

ในทางปฏิบัติ สัดส่วนการแบ่งข้อมูลที่นิยมใช้กันคือ 70:30 หรือ 80:20 ขึ้นอยู่กับขนาดของข้อมูลทั้งหมด หากข้อมูลมีขนาดใหญ่มาก อาจใช้สัดส่วน 90:10 ก็เพียงพอ แต่หากข้อมูลมีขนาดน้อย ต้องระวังไม่ให้ชุดทดสอบมีขนาดเล็กเกินไปจนทำให้ผลการประเมินไม่น่าเชื่อถือ นอกจากนี้ สิ่งสำคัญคือข้อมูลชุดทดสอบต้องไม่เคยถูกใช้ในการเรียนรู้ของตัวแบบก่อนหน้านี้เลย

\begin{myBox}[]{ขั้นตอนการทำ Train-Test Split}{}
\begin{enumerate}
\item สุ่มแบ่งข้อมูลทั้งหมดออกเป็น 2 ส่วน: ชุดฝึกสอน (Training Set) และชุดทดสอบ (Test Set)
\item ใช้ข้อมูลชุดฝึกสอนในการฝึก (Train) ตัวแบบ
\item ใช้ข้อมูลชุดทดสอบในการประเมิน (Evaluate) ประสิทธิภาพของตัวแบบที่ฝึกแล้ว
\item รายงานผลประสิทธิภาพที่ได้จากชุดทดสอบเป็นค่าประมาณของประสิทธิภาพที่แท้จริง
\end{enumerate}

สิ่งที่ควรระวังคือ ห้ามใช้ข้อมูลชุดทดสอบในการฝึกตัวแบบหรือปรับแต่งพารามิเตอร์ใด ๆ เพราะจะทำให้ผลการประเมินมีความลำเอียง (Biased) และไม่สะท้อนประสิทธิภาพที่แท้จริง
\end{myBox}

ข้อจำกัดของ Train-Test Split คือผลการประเมินขึ้นอยู่กับการสุ่มแบ่งข้อมูล หากการสุ่มแบ่งไม่ดี อาจทำให้ผลการประเมินไม่แม่นยำ นอกจากนี้ วิธีนี้ยังไม่ได้ใช้ข้อมูลทั้งหมดในการฝึกตัวแบบ ทำให้ตัวแบบอาจมีประสิทธิภาพลดลงเล็กน้อยเมื่อนำไปใช้กับข้อมูลทั้งหมด

\subsection{Cross-Validation}
Cross-Validation เป็นวิธีการตรวจสอบความถูกต้องที่มีประสิทธิภาพมากกว่า Train-Test Split โดยใช้ข้อมูลทั้งหมดในการฝึกและทดสอบตัวแบบหลายรอบ วิธีที่นิยมใช้กันมากที่สุดคือ K-Fold Cross-Validation

\begin{myBox}[]{K-Fold Cross-Validation}{}
ขั้นตอนของ K-Fold Cross-Validation มีดังนี้

\begin{enumerate}
\item แบ่งข้อมูลออกเป็น $K$ ส่วนเท่า ๆ กัน (เรียกว่า Folds)
\item ฝึกตัวแบบ $K$ รอบ โดยในแต่ละรอบจะใช้ข้อมูล $K-1$ Folds สำหรับฝึก และ 1 Fold สำหรับทดสอบ
\item คำนวณค่าประสิทธิภาพ (เช่น Accuracy, MSE) ในแต่ละรอบ
\item คำนวณค่าเฉลี่ยและส่วนเบี่ยงเบนมาตรฐานของค่าประสิทธิภาพทั้ง $K$ รอบ
\end{enumerate}

ค่า $K$ ที่นิยมใช้คือ 5 และ 10 ส่วนเบี่ยงเบนมาตรฐานของผลการทดสอบบอกความมั่นใจในผลการประเมิน ถ้าส่วนเบี่ยงเบนมาตรฐานต่ำ แสดงว่าผลการประเมินมีความคงที่ข้าม Fold ต่าง ๆ
\end{myBox}

ข้อดีของ Cross-Validation เปรียบเทียบกับ Train-Test Split คือ (1) ใช้ข้อมูลทั้งหมดในการทดสอบ ทำให้การประเมินมีความแม่นยำมากขึ้น (2) ลดผลกระทบจากการสุ่มแบ่งข้อมูลที่อาจเกิดความลำเอียง (3) ให้ข้อมูลเชิงสถิติเกี่ยวกับความแปรปรวนของประสิทธิภาพตัวแบบ

\begin{exmp}
พิจารณาข้อมูล 1000 ตัวอย่าง เมื่อใช้ 5-Fold Cross-Validation จะแบ่งข้อมูลออกเป็น 5 กลุ่ม กลุ่มละ 200 ตัวอย่าง การฝึกจะทำ 5 รอบ โดยในแต่ละรอบจะฝึกด้วยข้อมูล 800 ตัวอย่างและทดสอบด้วย 200 ตัวอย่างที่ไม่เคยเห็นในรอบนั้น สมมติผลการทดสอบในแต่ละรอบได้ค่า Accuracy ดังนี้: 0.85, 0.87, 0.83, 0.86, 0.84

ค่าเฉลี่ย Accuracy = (0.85 + 0.87 + 0.83 + 0.86 + 0.84)/5 = 4.25/5 = 0.85

ส่วนเบี่ยงเบนมาตรฐาน = $\sqrt{\frac{(0.85-0.85)^2+(0.87-0.85)^2+(0.83-0.85)^2+(0.86-0.85)^2+(0.84-0.85)^2}{5-1}} = \sqrt{0.0002} \approx 0.014$

ดังนั้นค่าประสิทธิภาพของตัวแบบโดยประมาณเท่ากับ 0.85 ± 0.014 ซึ่งแสดงความมั่นใจในผลการประเมินว่ามีความคงที่ดี
\end{exmp}

\begin{figure}[h]
\begin{pycode}{cross\_validation.py}
from sklearn.model_selection import (
    cross_val_score, KFold
)
from sklearn.ensemble import RandomForestClassifier
from sklearn.datasets import make_classification

X, y = make_classification(
    n_samples=1000, n_features=20,
    random_state=42
)

model = RandomForestClassifier(
    n_estimators=100, random_state=42
)

kf = KFold(n_splits=5, shuffle=True,
           random_state=42)

scores = cross_val_score(
    model, X, y, cv=kf, scoring='accuracy'
)

print(f"Accuracy per fold: {scores}")
print(f"Mean: {scores.mean():.4f}")
print(f"Std: {scores.std():.4f}")
\tcblower
ผลลัพธ์ที่ได้ (ตัวอย่าง):

Accuracy per fold: [0.85 0.87 0.83 0.86 0.84]
Mean: 0.8500
Std: 0.0141
\end{pycode}
\caption{Python code และผลลัพธ์การทำ Cross-Validation}
\label{fig:cv_code}
\end{figure}

%##########################################
\section{ปัญหา Overfitting และ Underfitting}
Overfitting และ Underfitting เป็นปัญหาพื้นฐานที่สำคัญที่สุดในการเรียนรู้ของเครื่อง ซึ่งเกี่ยวข้องโดยตรงกับความสามารถในการสังเคราะห์ (Generalization) ของตัวแบบ การเข้าใจธรรมชาติของปัญหาทั้งสองและสาเหตุที่ทำให้เกิดขึ้นเป็นสิ่งจำเป็นสำหรับการสร้างตัวแบบที่มีประสิทธิภาพดี

\begin{myLimit}[]{Overfitting และ Underfitting}{}
\begin{itemize}
\item \textbf{Overfitting} คือสถานการณ์ที่ตัวแบบ ``จำ'' ข้อมูลฝึกสอนได้ดีมากจนเกินไป รวมถึง Noise และความผันแปรสุ่มในข้อมูลฝึกสอน ทำให้ประสิทธิภาพกับข้อมูลฝึกสอนสูงมาก แต่ประสิทธิภาพกับข้อมูลใหม่ (Test Data) ต่ำกว่าที่คาดหวัง ในทางคณิตศาสตร์ Bias ต่ำแต่ Variance สูง

\item \textbf{Underfitting} คือสถานการณ์ที่ตัวแบบไม่ซับซ้อนพอที่จะจับรูปแบบที่แท้จริง (Pattern) ในข้อมูล ทำให้ประสิทธิภาพต่ำทั้งกับข้อมูลฝึกสอนและข้อมูลทดสอบ ในทางคณิตศาสตร์ Bias สูงแต่ Variance ต่ำ
\end{itemize}
\end{myLimit}

แนวคิด Bias-Variance Trade-off เป็นกรอบการวิเคราะห์ที่ช่วยให้เข้าใจปัญหานี้อย่างเป็นระบบ โดย Bias วัดความผิดพลาดโดยเฉลี่ยของตัวแบบจากค่าจริง ส่วน Variance วัดความไวของตัวแบบต่อการเปลี่ยนแปลงของข้อมูลฝึกสอน ตัวแบบที่ซับซ้อนมากเกินไปมักมี Variance สูง (Overfitting) ขณะที่ตัวแบบที่เรียบง่ายเกินไปมักมี Bias สูง (Underfitting) ตัวแบบที่ดีต้องมีสมดุลระหว่าง Bias และ Variance

สาเหตุที่ทำให้เกิด Overfitting ได้แก่ ข้อมูลฝึกสอนมีขนาดน้อยเกินไป ตัวแบบมีความซับซ้อนมากเกินไป (พารามิเตอร์มากเกินไป) หรือข้อมูลมี Noise สูง ส่วนสาเหตุของ Underfitting ได้แก่ ตัวแบบไม่ซับซ้อนพอ ข้อมูลไม่เพียงพอ หรือลักษณะสำคัญ (Features) ไม่เหมาะสม

วิธีแก้ปัญหา Overfitting ได้แก่ การเพิ่มข้อมูลฝึกสอน การลดความซับซ้อนของตัวแบบ การใช้เทคนิค Regularization การใช้ Cross-Validation ในการเลือกตัวแบบ และการตัดแต่ง (Pruning) สำหรับต้นไม้ตัดสินใจ ส่วนวิธีแก้ปัญหา Underfitting ได้แก่ การเพิ่มความซับซ้อนของตัวแบบ การเพิ่มจำนวนรอบการเรียน (Epochs) การใช้ Features ที่ดีขึ้น หรือการลดข้อจำกัดในตัวแบบ

%##########################################
\section{เทคนิค Regularization: Ridge และ Lasso}
Regularization เป็นเทคนิคที่ใช้ป้องกันปัญหา Overfitting โดยการเพิ่มข้อจำกัด (Constraint) ให้กับตัวแบบระหว่างการเรียนรู้ เป้าหมายคือการลดความซับซ้อนของตัวแบบโดยการลดค่าน้ำหนัก (Weights) ของพารามิเตอร์ที่ไม่จำเป็นให้ใกล้ศูนย์หรือเท่ากับศูนย์ เทคนิค Regularization ที่นิยมใช้กันมี 2 ชนิดหลักคือ Ridge และ Lasso

\begin{myBox}[]{Ridge Regression และ Lasso Regression}{}
สมมติให้มีฟังก์ชันสูญเสีย (Loss Function) สำหรับการถดถอยเชิงเส้นดั้งเดิมคือ

$$L(\boldsymbol{\beta}) = \sum_{i=1}^{n}(y_i - \beta_0 - \beta_1 x_{i1} - \cdots - \beta_p x_{ip})^2$$

\begin{enumerate}
\item \textbf{Ridge Regression} เพิ่มพจน์ L2 Regularization
$$L_{Ridge}(\boldsymbol{\beta}) = \sum_{i=1}^{n}(y_i - \beta_0 - \cdots - \beta_p x_{ip})^2 + \lambda\sum_{j=1}^{p}\beta_j^2$$

โดยที่ $\lambda \geq 0$ คือพารามิเตอร์ Regularization (ต้องปรับแต่ง) และ $\sum_{j=1}^{p}\beta_j^2$ คือ L2 Norm ของเวกเตอร์สัมประสิทธิ์ Ridge จะลดค่า $\beta_j$ ให้ใกล้ศูนย์แต่ไม่ทำให้เป็นศูนย์เด็ดขาด

\item \textbf{Lasso Regression} เพิ่มพจน์ L1 Regularization
$$L_{Lasso}(\boldsymbol{\beta}) = \sum_{i=1}^{n}(y_i - \beta_0 - \cdots - \beta_p x_{ip})^2 + \lambda\sum_{j=1}^{p}|\beta_j|$$

โดยที่ $\sum_{j=1}^{p}|\beta_j|$ คือ L1 Norm ของเวกเตอร์สัมประสิทธิ์ Lasso มีคุณสมบัติพิเศษคือสามารถทำให้บาง $\beta_j$ กลายเป็นศูนย์ได้ จึงเป็นทั้งวิธี Regularization และ Feature Selection ในเวลาเดียวกัน
\end{enumerate}
\end{myBox}

ความแตกต่างสำคัญระหว่าง Ridge และ Lasso อยู่ที่ผลกระทบต่อสัมประสิทธิ์ ใน Ridge Regression ค่าสัมประสิทธิ์ทั้งหมดจะถูกลดขนาดให้เล็กลงแต่ไม่มีสักตัวที่เป็นศูนย์อย่างแท้จริง ทำให้ตัวแบบยังคงใช้ตัวแปรทั้งหมด ในทางกลับกัน Lasso สามารถทำให้ค่าสัมประสิทธิ์บางตัวเป็นศูนย์ได้โดยตรง ซึ่งเหมาะสำหรับปัญหาที่มีตัวแปรอธิบายจำนวนมากและต้องการเลือกเฉพาะตัวแปรที่มีความสำคัญจริง ๆ ในทางเรขาคณิต Ridge ทำให้ค่าสัมประสิทธิ์ถูก ``ดึง'' เข้าหาจุดกำเนิดในรูปวงกลม ในขณะที่ Lasso ดึงเข้าหาจุดที่อยู่บนแกนทำให้บางมิติกลายเป็นศูนย์

\begin{exmp}
พิจารณาข้อมูลการถดถอยเชิงเส้นพหุคูณที่มีตัวแปรอธิบาย 10 ตัว รวมถึงตัวแปรที่ไม่มีความสัมพันธ์ที่แท้จริงกับตัวแปรตาม เมื่อใช้ Ridge Regression ด้วยค่า $\lambda$ ที่เหมาะสม พบว่าค่าสัมประสิทธิ์ของตัวแปรที่ไม่เกี่ยวข้องถูกลดลงให้ใกล้ศูนย์มาก แต่ยังคงมีค่าเล็กน้อย ส่วนตัวแปรที่มีความสัมพันธ์จริงยังคงมีค่าสัมประสิทธิ์ที่มีนัยสำคัญ เมื่อเปรียบเทียบกับ Lasso ค่าสัมประสิทธิ์ของตัวแปรที่ไม่เกี่ยวข้องถูกทำให้เป็นศูนย์โดยตรง ทำให้ตัวแบบมีเพียง 5 ตัวแปรที่ใช้งานจริง
\end{exmp}

\begin{figure}[h]
\begin{pycode}{regularization.py}
import numpy as np
from sklearn.linear_model import Ridge, Lasso
from sklearn.model_selection import cross_val_score
from sklearn.preprocessing import StandardScaler

# Standardize features for regularization
scaler = StandardScaler()
X_scaled = scaler.fit_transform(X)

# Ridge Regression with different alpha values
alphas = [0.01, 0.1, 1, 10, 100]
for alpha in alphas:
    ridge = Ridge(alpha=alpha)
    scores = cross_val_score(
        ridge, X_scaled, y, cv=5,
        scoring='neg_mean_squared_error'
    )
    rmse = np.sqrt(-scores.mean())
    print(f"Ridge alpha={alpha}: RMSE={rmse:.4f}")

# Lasso Regression
for alpha in alphas:
    lasso = Lasso(alpha=alpha)
    scores = cross_val_score(
        lasso, X_scaled, y, cv=5,
        scoring='neg_mean_squared_error'
    )
    rmse = np.sqrt(-scores.mean())
    n_features = np.sum(lasso.coef_ != 0)
    print(f"Lasso alpha={alpha}: RMSE={rmse:.4f}, "
          f"n_features={n_features}")
\tcblower
ผลลัพธ์ที่ได้ (ตัวอย่าง):

Ridge alpha=0.01: RMSE=5.2341
Ridge alpha=0.1: RMSE=4.8923
Ridge alpha=1: RMSE=4.1234
Ridge alpha=10: RMSE=4.5678
Ridge alpha=100: RMSE=5.8901

Lasso alpha=0.01: RMSE=4.0892, n\_features=9
Lasso alpha=0.1: RMSE=4.2345, n\_features=7
Lasso alpha=1: RMSE=4.6789, n\_features=5
Lasso alpha=10: RMSE=5.4321, n\_features=2
Lasso alpha=100: RMSE=6.1234, n\_features=0

หมายเหตุ: n\_features คือจำนวนตัวแปร
ที่มีสัมประสิทธิ์ไม่เท่ากับศูนย์
\end{pycode}
\caption{Python code และผลลัพธ์การเปรียบเทียบ Ridge และ Lasso}
\label{fig:regularization_code}
\end{figure}

%##########################################
\section{การเปรียบเทียบและคัดเลือกตัวแบบ}
ในทางปฏิบัติ การพัฒนาตัวแบบการเรียนรู้ของเครื่องมักไม่ได้มีเพียงตัวเลือกเดียว แต่ต้องพิจารณาตัวแบบหลายตัวที่มีประสิทธิภาพแตกต่างกัน การเปรียบเทียบและคัดเลือกตัวแบบที่เหมาะสมจึงเป็นขั้นตอนสำคัญในกระบวนการวิทยาการข้อมูล

หลักการในการเปรียบเทียบตัวแบบมีดังนี้ ประการแรก ต้องใช้ข้อมูลชุดเดียวกันในการทดสอบทุกตัวแบบเพื่อให้การเปรียบเทียบมีความยุติธรรม ประการที่สอง ควรใช้ตัวชี้วัดประสิทธิภาพที่เหมาะสมกับประเภทปัญหา เช่น Accuracy หรือ F1-Score สำหรับ Classification และ MSE หรือ R-squared สำหรับ Regression ประการที่สาม ต้องใช้การตรวจสอบความถูกต้อง (Validation) เช่น Cross-Validation เพื่อให้ได้ผลการประเมินที่น่าเชื่อถือ และประการสุดท้าย ต้องพิจารณาปัจจัยอื่น ๆ นอกเหนือจากประสิทธิภาพ เช่น ความเรียบง่ายในการตีความ (Interpretability) เวลาในการฝึกและทำนาย รวมถึงความต้องการทรัพยากรคอมพิวเตอร์

วิธีการคัดเลือกตัวแบบที่นิยมใช้ ได้แก่ การเปรียบเทียบโดยตรงจากผล Cross-Validation การใช้เกณฑ์ BIC (Bayesian Information Criterion) หรือ AIC (Akaike Information Criterion) สำหรับตัวแบบทางสถิติ การใช้ Ensemble Methods เช่น Voting หรือ Stacking เพื่อรวมตัวแบบหลายตัว และการใช้เทคนิค Hyperparameter Tuning เช่น Grid Search หรือ Random Search เพื่อหาค่าพารามิเตอร์ที่เหมาะสมที่สุด

เมื่อเลือกตัวแบบที่ดีที่สุดได้แล้ว ขั้นตอนสุดท้ายคือการนำตัวแบบไปฝึกด้วยข้อมูลทั้งหมด (Training Set + Test Set) ก่อนนำไปใช้งานจริง เนื่องจากในการประเมินประสิทธิภาพที่ผ่านมา เราถือส่วนหนึ่งของข้อมูลไว้สำหรับทดสอบ ทำให้ตัวแบบที่ได้ยังไม่ได้ใช้ข้อมูลทั้งหมดในการเรียนรู้ ดังนั้นเมื่อได้ตัวแบบที่ดีที่สุดแล้ว ควรฝึกใหม่ด้วยข้อมูลทั้งหมดเพื่อให้ได้ตัวแบบขั้นสุดท้ายที่พร้อมใช้งาน

\begin{exercise}
\begin{enumerate}
\item จากข้อมูลการพยากรณ์ราคาบ้าน 10 ตัวอย่าง จงคำนวณค่า MSE, RMSE, MAE และ R-squared เมื่อค่าจริงคือ [250, 300, 180, 420, 380, 200, 450, 320, 280, 390] (หน่วย: พันบาท) และค่าทำนายคือ [245, 310, 175, 400, 370, 210, 430, 315, 290, 385]

\item จงอธิบายความแตกต่างระหว่าง Overfitting และ Underfitting พร้อมยกตัวอย่างสถานการณ์ที่อาจเกิดปัญหาแต่ละประเภท

\item ทำไม F1-Score จึงเป็นตัวชี้วัดที่ดีกว่า Accuracy ในบางสถานการณ์ และจงยกตัวอย่างสถานการณ์ที่ Accuracy ไม่เหมาะสมที่จะใช้

\item เปรียบเทียบข้อดีข้อดีของ Ridge และ Lasso Regression โดยในสถานการณ์ใดควรเลือกใช้แต่ละวิธี
\end{enumerate}
\end{exercise}

\newpage
%\RefChapter
